% PART 12 DRY RUN INPUT. Not the Overleaf file itself.
% Source: the paper source as the authors had it at 2026-09-23T05:58:23Z (22:58 PDT, 22 Sep),
% copied from the authors' working notes, which are not released.
% It was the only full copy of the source (\documentclass to \end{document}) in those notes.
% The dry run tested the Part 12 scripts on it before the final tex was saved.
% Characters 895..41286 of the copied text were kept. Text before \documentclass
% and after \end{document} dropped. Body is byte-for-byte as copied.
% sha256(body) = 8fbc0115606cda474f48d32c392caf6331993f35b02a9958feb26adaace64c12
% Line numbers in Part 12 dry-run outputs refer to THIS file (body starts at line 10).
\documentclass{article}
\usepackage{spconf,amsmath,amssymb,graphicx,booktabs,microtype,url,comment,enumitem,xcolor}

\graphicspath{{figures/}}
% tighter float and equation spacing to keep the text on four pages
\setlength{\textfloatsep}{8pt plus 2pt minus 2pt}
\setlength{\abovecaptionskip}{4pt}
\setlength{\belowcaptionskip}{2pt}
\makeatletter
\g@addto@macro\normalsize{\setlength{\abovedisplayskip}{4pt}\setlength{\belowdisplayskip}{4pt}\setlength{\abovedisplayshortskip}{2pt}\setlength{\belowdisplayshortskip}{2pt}}
\makeatother

\begin{document}

\title{The Representation Utility Gap in Audio LLMs for Clinical Speech}

\name{Chaitanya Parwatkar, Nima Kelidari, Minoo Ahmadi, Ashutosh Chaubey, Mohammad Soleymani%
  \thanks{Corresponding author: soleymani@ict.usc.edu.}}
\address{University of Southern California, Los Angeles, CA, USA}

\maketitle

\begin{abstract}
%Recent advances in speech processing have demonstrated the utility of speech as a robust marker of clinical conditions that affect speech production. Audio LLMs can answer clinical questions about a recording in natural language. However, Audio LLMs' output depends on both paralinguistic and linguistic cues, and it is unclear how much of their output is grounded in paralinguistic rather than linguistic cues. In this paper, we compare the information present in audio representations against the responses generated by the LLMs. We train linear probes on the frozen hidden states of an Audio LLM and query models across three conditions that differ in their language saliency: Parkinson's read-speech, depression interviews and Alzheimer's picture descriptions. Probes detect disorders from audio representations with AUC ranging from 0.75 to 0.92, whereas the LLM's AUC is much lower due to a representation utility gap. Where the words contradict the correct diagnosis, LLMs completely fail. To narrow this utility gap, we fine-tune the projector layer, thereby boosting performance. This points to a utility gap between the information in speech representations and the outputs of Audio LLMs.%
Speech is an established marker of clinical conditions that affect speech production, and Audio LLMs can now answer clinical questions about a recording in natural language. Their output depends on both paralinguistic and linguistic cues, and it is unclear how much of it is grounded in the voice rather than the words. We compare the information present in an Audio LLM's audio representations against the answers it generates. We train linear probes on its frozen hidden states and query it under three conditions that differ in language saliency: Parkinson's read-speech, depression interviews and Alzheimer's picture descriptions. Probes detect the disorders from the representations with AUC from 0.75 to 0.92, while the zero-shot answer sits well below on Parkinson's and Alzheimer's, a representation utility gap. Where the words contradict the diagnosis, four of six models fall below chance. Retraining only the projector narrows the gap on Parkinson's and Alzheimer's.
\end{abstract}

\begin{keywords}
audio language models, clinical speech,  Alzheimer's disease, Parkinson's disease, depression
\end{keywords}

\section{Introduction}\label{sec:intro}

Recent advances in speech processing demonstrated the utility of speech as a robust marker for a range of disorders, including psychological distress \cite{Cummins2015} and neurodegenerative conditions \cite{9426437}. Audio language models (Audio LLMs) now answer clinical questions about a recording in natural language and can reason about their decisions \cite{loweimi_zero-shot_2025}. The ability of large foundation models to perform zero-shot detection and explain their reasoning makes them a potentially attractive tool for clinical applications. However, the validity of reasoning and grounding in auditory evidence has not been established. 

\begin{figure}[!t]
\centering
\includegraphics[width=\columnwidth]{fig1_gap.pdf}
\caption{\small Linear probe AUC at each stage of Qwen2.5-Omni, speaker-disjoint folds, one dataset per condition. The 32 audio encoder stages come first, then the projector, then the 29 language model stages read at the audio-token positions.}
\label{fig:gap}
\end{figure}

Current AI-based diagnosis approaches are primarily evaluated on their accuracy and little attention is given to what the models use to reach their decisions. Audio LLMs in particular are able to leverage both auditory and linguistic information from speech as they are very well capable of transcribing speech, even when none is given as prompt. Hence, it is not evident if their decisions are grounded in paralinguistic rather than linguistic evidence, when spoken words and paralinguistic speech can both point towards the presence of a clinical condition.

There is a wide body of research in deep speech-based models for clinical assessment \cite{Cummins2015}. Recently, large pre-trained models and multimodal LLMs have become the dominant tools for AI-based clinical assessment \cite{speechdx,fend,loweimi_zero-shot_2025}. Bhalla et al. \cite{speechdx} have developed a large benchmark and evaluated a set of large pretrained audio encoders for the diagnosis of clinical conditions, including neuro-degenerative disorders that disrupt speech production. They found that no single audio encoder generalizes across the clinical speech landscape and urge further research on developing a generalizable audio encoder for health applications. Dong et al. \cite{fend} created a large multi-lingual benchmark for Alzheimer's disease (AD) detection and similarly evaluated the performance of a deep speech encoder for their diagnosis through multimodal fusion of speech and language. They found that fusion fails to beat unimodal models, and the dominance of language or paralinguistic speech is domain-dependent. These benchmarks were used to measure how well a representation encodes a condition, not whether the model's answer uses it. Chen et al. \cite{listen}, and Pang et al. \cite{voxparadox} show that Audio LLMs follow the words in paralinguistic tasks, and the latter attribute this to signal degradation in the deep encoder layers. We extend their study to clinical diagnosis and evaluate conflicting linguistic and paralinguistic signals that arise naturally in speech.

In this paper, we study how Audio LLMs leverage information in audio representations to arrive at an answer about a clinical diagnosis. Linear probes are trained on the frozen hidden states of an Audio LLM, which is then queried on the same recordings with the same question. We perform these comparisons under three different conditions so that the diagnostic content of the spoken words differs between them, including Parkinson's read-speech, where all speakers read the same passage and the spoken words are irrelevant; depression interviews, where the words carry part of the information necessary for diagnosis; and Alzheimer's picture descriptions, where the words carry most of the necessary information. %We find that the paralinguistic signal is present in the audio representations but is not reflected in the answer, which instead follows the spoken words. We trace this utility gap to the final step that maps the model's states onto a language token of the answer. We then demonstrate that the gap can be narrowed by fine-tuning a single component of the model, the projector that sits between the audio encoder and the LLM.  
The major contributions of this work are as follows.
\begin{enumerate}[leftmargin=*,itemsep=0pt,topsep=0pt,parsep=0pt]
\item We construct a set of samples with conflicting audio and language content from real clinical recordings that allows a systematic study of Audio LLMs' reliance on linguistic vs paralinguistic speech.
\item We compare what an Audio LLM encodes against what it answers on seven clinical datasets, and test six Audio LLMs on the conflict sets to test whether the paralinguistic signal is utilized.
\item We localize the utility gap to the final readout of the LLM from representations and show that retraining only the projector recovers most of the performance.
\end{enumerate}


\section{Method}\label{sec:setup}

We aim to establish whether an Audio LLM's answer to speech-based clinical diagnosis is grounded in the auditory evidence its representations hold. To this end, we first read the diagnostic information from the frozen model's hidden states with a linear probe, trained and evaluated on speaker-disjoint folds probing every layer of the encoder and the language model. For layer $\ell$ we mean-pool the hidden states over time and fit a logistic probe on speaker-disjoint folds and the AUC is calculated as follows.
\begin{equation}
\mathrm{AUC}^{\text{probe}}_\ell = \mathrm{AUC}\big(\sigma(\mathbf{w}_\ell^{\top}\bar{\mathbf{h}}_\ell + b_\ell),\, y\big)
\label{eq:probe} \end{equation}
where $y$ is the  label and $\bar{\mathbf{h}}_\ell$ is the mean-pooled activation. We then ask the model the same prompt (clinical question) and read its answer from the logits at the first answer position, in one forward pass and with no generation, 
\begin{equation}
\mathrm{AUC}^{\text{answer}} = \mathrm{AUC}(p(\text{Y}),\, y), \:\: p(\text{Y}) = \sigma\big(z(\text{Y}) - z(\text{N})\big) \label{eq:answer} 
\end{equation} 
$z(\text{Y})$ and $z(\text{N})$ are the logits of the first token of \emph{Yes} and \emph{No} at the first answer position, so $p(\text{Y})$ is the probability of \emph{Yes} conditioned on the answer being one of the two.  We also train a probe on the hidden state at this position. The model reads its Yes or No answer from that vector.

The utility gap is defined as $\Delta_\ell = \mathrm{AUC}^{\text{probe}}_\ell - \mathrm{AUC}^{\text{answer}}$

In both cases we use frozen models, identical clips and an identical prompt, so any difference between the results is a property of the model rather than of the data or the prompt.

We repeat the comparison across six Audio LLMs that share this encoder, projector and LLM design, to establish whether the utility gap is common across architectures. The models are Qwen2-Audio \cite{qwen2audio}, Qwen2.5-Omni \cite{qwen25omni}, Qwen3-Omni \cite{qwen3omni}, Audio Flamingo 2 and 3 \cite{af2,af3} and Kimi-Audio \cite{kimiaudio}.

Such a comparison on its own cannot evaluate whether a model relies on paralinguistic or linguistic speech, since the spoken words and the diagnosis usually point in the same direction in clinical recordings. Hence, we add two conditions that allow a conflict between acoustics and language. The conflict sets are created for depression and Alzheimer's disease. Each conflict segment is paired with an agreement segment from the same speaker. Given that the speech signals are real, the contradictions occur naturally. We also run a transcript only baseline where the model gets the words alone, with the same prompt and no audio, and establish how much of the decision can be attributed to language alone. Finally, we ask whether the gap closes when the projector alone is retrained (see Section~\ref{sec:results}).

\section{Implementation Details}\label{sec:impl}
\noindent\textbf{Datasets.} Seven datasets across Parkinson's disease (PD), Alzheimer's disease (AD) and major depressive disorder (MDD), are selected so that the diagnostic content of the spoken words differs between them. 

From the Parkinson's disease datasets, we only used the read-speech (read-aloud) protocol where every speaker reads the same passage, so the words carry almost no diagnostic signals and any signal must come from the voice. We use the three following corpora: PC-GITA (Spanish, 1100 clips) \cite{pcgita}, NeuroVoz (Spanish, 1270 clips) \cite{neurovoz} and MDVR-KCL (English, 37 speakers) \cite{kcl}.
 
The clinical interviews for depression diagnosis sit in between, with language containing information that points to the diagnosis. We use E-DAIC \cite{daic}, the extended corpus released for AVEC 2019, with 275 participants, where the binary label is the dataset's PHQ-8 \cite{Kroenke2009} based label rather than a clinical diagnosis.  The model hears the participant's whole interview, capped at 15 minutes for the 89 longest.
 
For Alzheimer's disease we use picture description, where semantic content and fluency carry most of the diagnosis. We use the DementiaBank Pitt cookie-theft descriptions \cite{pitt}, with 468 segments from 228 speakers, the top and bottom third of the fluency score described below, the ADReSSo recordings \cite{adresso} from 237 speakers, as released with the ADReSS-M challenge \cite{adressm}, and the ADReSS-2020 benchmark split \cite{adress2020} including 108 training and 48 test recordings. Each Pitt segment starts at the participant's first timed utterance, and the ADReSSo and ADReSS-2020 windows start at the beginning of the recording, so interviewer prompts inside the window are kept. Models hear the first 30\,s, the length Qwen2-Audio accepts, so the conflict test in Table~\ref{tab:conflict} runs on the same clips for all six models. The three sets share speakers, so every result is within one set on speaker-disjoint folds and we make no transfer claim among them.
 
%An eighth corpus, an Italian Parkinson's dataset, was excluded because recording properties alone separated its groups at 0.995 (Section~\ref{sec:controls}).
 
\noindent\textbf{Probing.} The probe of Eq.~\ref{eq:probe} is logistic regression with $\ell_2$ penalty at $C=1$, balanced class weights and 2000 iterations, on features standardized inside the training fold. Folds are speaker-disjoint (GroupKFold, 5 folds) and the layer is chosen inside the training folds. %Where a layer is chosen, it is chosen inside the training folds; selecting the layer after seeing the test fold inflates AUC by $+0.03$ on average across 18 cells and by up to $+0.14$, so we quote the nested value wherever both exist.
The layer-wise probing is done on Qwen2.5-Omni \cite{qwen25omni}, which we probe end-to-end: 32 hidden states in its audio encoder, the projector, and 29 hidden states in its language model taken at the audio-token positions. We chose this model due to its simpler architecture allowing for consistent probing, since Qwen3-Omni \cite{qwen3omni} uses Mixture-of-Experts.

\noindent\textbf{Fine-tuning.} The projector is trained with an AdamW optimizer with a learning rate $10^{-4}$, for three epochs, one clip per update and cross entropy over the ``Yes'' and ``No'' logits only, on the same 5 speaker-disjoint folds.

\noindent\textbf{Prompts.} \emph{``Based only on this recording, does this speaker show signs of X? Answer with one word, Yes or No''}; on the Parkinson's and depression datasets, the phrase ``this recording'' is replaced with ``how this person's voice sounds''.  For the transcript only baseline, ``this recording'' is replaced with ``this transcript of what the person said''. Two other wordings move the answers by 0.02 on average. %The two score within 0.01 of each other on Pitt.


\noindent\textbf{Conflict set construction.} From the raw E-DAIC transcripts, with the interviewer removed, we take every participant segment of at least 25 words and 5\,s of speech and score it with a RoBERTa sentiment classifier trained on TweetEval \cite{tweeteval}. A segment counts as positive or negative when that class has a probability of at least 0.70. A participant with PHQ-8 of 15 or more speaking positively, or a healthy control with PHQ-8 of 4 or less speaking negatively, is a conflict segment. This yields 208 of the first kind and 280 of the second, each paired with an agreement segment from the same speaker: 483 pairs over 138 speakers, so that across the whole set the sentiment of the words does not predict the label. For Alzheimer's disease, we score the words of each segment in Pitt database \cite{pitt} with a logistic regression trained on six fluency markers from the CHAT transcription codes \cite{macwhinney2000chat,fraser2016linguistic}: type token ratio, retraces, fillers, unintelligible tokens, errors and pauses, with speaking rate excluded because it alone reaches 0.92 and would make the split trivial. The markers are standardized and the regression is trained on five speaker-disjoint folds, so each segment is scored out of fold. Segments whose out-of-fold score falls in the top or bottom third are kept. A participant with dementia scored as fluent, or a healthy control scored as impaired, is a conflict segment, and the reverse is an agreement segment. This leaves 146 segments in conflict and 322 in agreement. The median segment fits inside the 30\,s window, and 91 percent of segments have at least half their words inside it.

\noindent\textbf{Transcripts.} For the Pitt database, ADReSS-2020 \cite{adress2020} and E-DAIC \cite{daic}, we use the transcripts provided with the dataset after stripping the annotation codes; ADReSSo \cite{adresso} and the three Parkinson's sets have no transcripts, so we transcribe them with whisper-large-v3 \cite{whisper}.

\noindent\textbf{Checks.} As an acoustic baseline we train a separate logistic regression classifier on eGeMAPS features \cite{egemaps}. On the Pitt dataset participants with a positive diagnosis are seven years older than healthy controls on average, so we repeat the probe and the answer on 275 clips from 129 age and sex matched speakers.

\section{Experimental Results}\label{sec:results}

%\subsection{Controls on the measurement}\label{sec:controls}

%Before comparing probe against answer we check that the datasets support the comparison. Six recording measurements containing no speech content separate patients from controls at 0.995 on the Italian dataset; we exclude it. On the four Parkinson's and depression datasets the recording-only floor is 0.45 to 0.80 and the probe stays above it; on Neurovoz, age matching moves the encoder probe from 0.95 to 0.91 while a silence-only probe falls from 0.72 to 0.58. On Pitt the six together reach 0.65, below the probe (0.78); controls are younger (64 against 72 years, age alone 0.70), but on 64 age and sex matched speaker pairs (268 segments) the probe still reads 0.79 and the answer 0.66. A classical baseline, eGeMAPS features \cite{egemaps} into the same logistic classifier, reads 0.62 on Pitt, 0.73 on ADReSSo, 0.70 on ADReSS-2020 and 0.80 on MDVR-KCL; the probes beat it on every set, and on Pitt the model's own answer merely matches it.

\begin{table}[t]
\centering\small
\setlength{\tabcolsep}{5pt}
\caption{\small AUC of Qwen2.5-Omni per dataset. Zero-shot audio and text take the recording and its transcript. Fine-tuned is the same model after retraining only its projector on speaker-disjoint folds.}
\label{tab:main}
\begin{tabular*}{\columnwidth}{@{\extracolsep{\fill}}lrccc@{}}
\toprule
 & & \multicolumn{2}{c}{Zero-shot} & Fine-tuned \\
\cmidrule(lr){3-4}\cmidrule(l){5-5} Dataset & Clips (speakers) & Audio & Text & Audio \\
\midrule
\multicolumn{5}{@{}l}{\emph{Parkinson's disease}} \\
PC-GITA     & 1100 (100) & 0.56 & 0.52 & 0.89 \\
NeuroVoz    & 1270 (107) & 0.70 & 0.54 & 0.90 \\
MDVR-KCL    &   37 (37)  & 0.71 & 0.76 & 0.70 \\
\addlinespace[3pt]
\multicolumn{5}{@{}l}{\emph{Major depressive disorder}} \\
E-DAIC      &  275 (275) & 0.83 & 0.82 & 0.64 \\
\addlinespace[3pt]
\multicolumn{5}{@{}l}{\emph{Alzheimer's disease}} \\
Pitt        &  468 (228) & 0.66 & 0.66 & 0.77 \\
ADReSSo     &  237 (237) & 0.62 & 0.55 & 0.82 \\
ADReSS-2020 &  156 (156) & 0.62 & 0.75 & 0.77 \\
\bottomrule
\end{tabular*}
\end{table}

\subsection{Audio LLM performance}
We prompted Audio LLMs with a text prompt (clinical question) and a speech excerpt (clip) to perform zero-shot diagnosis. Zero-shot, Qwen2.5-Omni stays between 0.56 and 0.71 on the Parkinson's and Alzheimer's sets, and only E-DAIC is higher (Table~\ref{tab:main}). 

Table~\ref{tab:main} also reports the answer from the transcript alone. For AD the transcript alone matches the audio on the Pitt dataset, beats it on ADReSS-2020 and trails it on ADReSSo (Table~\ref{tab:main}). On the Parkinson's sets the transcript sits near chance, since every speaker reads the same sentences. On E-DAIC the audio answer on the whole interview matches the transcript alone. Qwen3-Omni and Audio Flamingo 3 reach 0.87 and 0.86 on the same interviews, while Qwen2-Audio, which hears 30\,s, stays at 0.63.


%Their performance also improves with having access to longer audio tracks. With 300\,s segments on depression, Audio Flamingo 3, Qwen2.5-Omni and Kimi-Audio achieve AUC of $0.86$, $0.84$ and $0.83$, respectively. The audio answer beats the transcript answer only where there are few words to read, on Parkinson's read-speech, where the transcript answer sits between 0.47 and 0.67; on depression no model beats its own transcript.
%Agreed, cutting it to two sentences: the recording-only floor and the age-matched check. -- MSL please add your comment to the comments not in the textt

\subsection{Probing Results}\label{sec:gap}
%The encoder holds what the answer does not use (Figure~\ref{fig:gap}). Under nested layer selection the encoder probe reads 0.92 on PC-GITA, 0.96 on Neurovoz and 0.89 on MDVR-KCL read-speech, and 0.78 on the Pitt Alzheimer's segments; the same model's answers on the same clips read 0.54, 0.55, 0.66 and 0.62. On Alzheimer's the signal does not stop at the encoder: probing the language model's own stages at the audio-token positions, the  readable signal rises to 0.86 by the deep layers while the answer stays where it was. On Alzheimer's the words also carry the task (a text classifier reads 0.70) and the model can transcribe, so the words could be helping here too; where the words carry almost nothing, on MDVR-KCL read-speech, the probe still holds 0.82 to 0.88 through every language-model stage. On depression at the 30\,s window Qwen2-Audio accepts, the nested probe (0.66) and the answer (0.65) coincide.

The best probe achieves an AUC of 0.75 to 0.92 on the Parkinson's sets and 0.77 on the Pitt dataset, far above the model's own answers on the same clips (Fig.~\ref{fig:gap}).  On all 708 Pitt segments, middle third included, the encoder probe reaches 0.80 and the zero-shot answer 0.67. The same probes on Qwen2-Audio achieve 0.77 to 0.93 against answers of 0.54 to 0.68 on the Parkinson's and Alzheimer's sets. On Qwen3-Omni they obtain AUC scores of 0.77 to 0.94 against 0.60 to 0.78. We note that the probe is trained in-domain; hence, comparing its performance to zero-shot recognition by Audio LLMs is not fair. However, these results demonstrate that representations contain necessary information for diagnosis. The encoder probe beats the eGeMAPS baseline on all seven datasets, by 0.06 on PC-GITA to 0.18 on ADReSSo.  Recording features from the quietest frames separate the NeuroVoz groups at 0.87 by spectral centroid alone and PC-GITA at 0.69, so part of the Parkinson's probe may come from the recording channel. The Alzheimer's and E-DAIC sets stay below 0.64 on every such feature. Age alone separates the Pitt groups at 0.71. On the age matched clips the age alone falls to 0.39, while the encoder probe, with its layer chosen inside the training folds, reaches 0.80 and the zero-shot answer 0.65. The information also survives past the audio encoder as probing the language model's representations at the audio-token positions on Alzheimer's rises to 0.80, and the answer state itself achieves 0.77. The probe and the answer use different directions of that vector. The cosine between the probe direction and the Yes minus No readout direction is 0.01 in magnitude, no larger than between two random directions. Removing the readout direction from the state leaves the probe at 0.77. The gap is widest on Parkinson's read-speech, narrower on Alzheimer's picture descriptions and reversed on depression. The performance difference shrinks as the words carry more diagnostic information. On E-DAIC the best probe is 0.69, at the language model, below the answer of 0.83.

%\subsection{Text only}

%The answer follows the words, and two controls agree. \emph{Transcript only.} Given the Pitt transcript and no audio, Qwen2-Audio scores 0.68 (95\% bootstrap interval 0.63 to 0.72); with the audio it scores 0.62 (0.56 to 0.67), and on the same clips the audio answer sits 0.06 below the transcript answer (paired interval $-0.12$ to $+0.01$): the audio adds nothing the words did not already carry. On ADReSSo the pair is 0.68 and 0.66. On E-DAIC the transcript-only answer of Qwen2.5-Omni correlates with its own audio answer at $r = 0.83$ and returns the identical Yes/No on 96\% of segments, and across the depression grid the median audio answer (0.83) sits beside the median transcript answer (0.81). \emph{Words check.} A probe trained only on Pitt segments where the words agree with the diagnosis reaches 0.95 within that set and at most 0.58 on the segments where the words contradict it; on ADReSSo, with one clip from each speaker, the same test gives 0.99 within and 0.39 to 0.61 on conflict, on a small conflict group of 19. What the probe learns here is largely the words, and they enter at the encoder itself, which was trained to transcribe.

\subsection{Conflict Test Results}

\begin{table}[t]
\centering\small
\setlength{\tabcolsep}{4pt}
\caption{\small AUC on segments where the words conflict with the diagnosis and on matched agreement segments from the same speakers, E-DAIC \cite{daic} for depression and Pitt \cite{pitt} for Alzheimer's.}
\label{tab:conflict}
\begin{tabular*}{\columnwidth}{@{\extracolsep{\fill}}lcccc@{}}
\toprule
 & \multicolumn{2}{c}{Depression} & \multicolumn{2}{c}{Alzheimer's} \\
\cmidrule(lr){2-3}\cmidrule(l){4-5}
Model & Conflict & Agree. & Conflict & Agree. \\
\midrule
Qwen2-Audio \cite{qwen2audio}   & 0.23 & 0.95 & 0.41 & 0.70 \\
Qwen2.5-Omni \cite{qwen25omni}  & 0.22 & 0.95 & 0.53 & 0.71 \\
Qwen3-Omni \cite{qwen3omni}     & 0.30 & 0.96 & 0.61 & 0.82 \\
Audio Flamingo 2 \cite{af2}     & 0.53 & 0.61 & 0.63 & 0.55 \\
Audio Flamingo 3 \cite{af3}     & 0.42 & 0.81 & 0.61 & 0.59 \\
Kimi-Audio \cite{kimiaudio}     & 0.53 & 0.87 & 0.72 & 0.72 \\
\bottomrule
\end{tabular*}
\end{table}

% Six of nine models fall to 0.36 or lower where the words point away from the label, as low as 0.29, while scoring 0.87 to 0.93 on the matched segments where words and label agree (Table~\ref{tab:conflict}); MiMo-Audio scores 0.47 and Kimi-Audio 0.52 on conflict. A model that listened would not invert; a model that reads must. On the Pitt dataset conflict segments Qwen2-Audio falls to 0.41 against 0.70 on agreement; qwen2.5-omni scores 0.41 against 0.75 on the same segments, answering No on every clip.
To test for the language bias in Audio LLMs, we tested their performance on a set of samples for MDD and AD detection where the spoken words do not point towards the diagnosis, such as positive sentiment for MDD and fluent speech for AD. A model attending to the paralinguistic speech should perform as it does on the matched agreement segments, whereas a model following the words should invert and fall below chance. No such test is possible for PD, since every speaker reads the same passage. Results in Table~\ref{tab:conflict} show a strong language bias. On the depression segments four of the six models fall below chance when the words point away from the label, while every model except Audio Flamingo 2 scores above 0.80 when the words and the label agree.  Speaker-level bootstrap intervals for the difference between the two arms exclude zero for all six models. Audio Flamingo 2 is the only model whose audio encoder, CLAP \cite{clap}, was not trained for speech recognition, so its transcription is limited. On the Pitt conflict segments the three Qwen models drop by 0.18 to 0.29 against their agreement segments, Kimi-Audio holds, and both Audio Flamingo models sit near chance on both sets (Table~\ref{tab:conflict}).


\begin{comment}

\begin{table}[t]
\centering\footnotesize
\caption{Conflict and agreement answers under each fix (AUC); the layer mixers, scored on PC-GITA
overall AUC, are in the text. Training-free fixes that move conflict pay elsewhere: subtraction and
steering in the agreement arm, the untrained plug in the overall answer (0.49). The omni baseline is
the sweep's rerun (0.31 in Table~\ref{tab:conflict}).}
\label{tab:fixes}
\setlength{\tabcolsep}{4pt}
\begin{tabular}{@{}lcc@{}}
\toprule
fix & conflict & agreement \\
\midrule
\multicolumn{3}{@{}l}{\emph{depression, no training, Qwen2.5-Omni}} \\
the model's own answer & 0.30 & 0.87 \\
best of 12 prompt wordings & 0.34 & 0.80 \\
text subtraction, $\alpha = 0.9$ & 0.50 & 0.72 \\
\multicolumn{3}{@{}l}{\emph{depression, no training, Qwen2-Audio}} \\
the model's own answer & 0.35 & 0.92 \\
best of 12 prompt wordings & 0.48 & 0.85 \\
activation steering, best strength & 0.54 & 0.50 \\
\midrule
\multicolumn{3}{@{}l}{\emph{Alzheimer's, no training, Qwen2-Audio}} \\
the model's own answer & 0.41 & 0.70 \\
few-shot audio examples (Yes on 1\%) & 0.46 & 0.45 \\
best encoder layer plugged in, untrained & 0.55 & 0.47 \\
steering at the probe's layer, $c = 1$ & 0.56 & 0.42 \\
\midrule
\multicolumn{3}{@{}l}{\emph{Alzheimer's, a small part trained, Qwen2-Audio}} \\
linear readout on the frozen states & \textbf{0.61} & 0.86 \\
projector retrained on layer 24 & \textbf{0.57} & 0.89 \\
\bottomrule
\end{tabular}
\end{table}
\end{comment}


%\subsection{Test-time Remedy Results}
%Seven training-free fixes fail (Table~\ref{tab:fixes}). Twelve prompt wordings never lift the depression conflict AUC above 0.48. Learned layer mixing over the encoder reaches 0.909 against 0.916 for one layer chosen inside the training folds, and conditioning the mixture on the prompt adds 0.0004. Subtracting the transcript-only logits from the audio logits stops the model following the words at weight $\alpha = 0.9$ (agreement with the lexical direction falls from 0.70 to 0.50), where conflict AUC is 0.50; larger $\alpha$ reaches 0.67 only by collapsing the agreement arm to 0.18, and held-out speakers choose $\alpha = 0$ in four of five folds. Adding a multiple of the difference-in-means diagnosis direction to the encoder output damages the agreement arm at the smallest coefficient before conflict moves, and a random direction of the same norm displaces the answer equally; across 20 steered conditions the apparent gain correlates at $-0.92$ with the lexical signal being destroyed, and on Pitt, at the probe's own layer 22, the diagnosis direction gives 0.56 on conflict and 0.42 on agreement against 0.58 and 0.54 for a random one. Feeding the probe's best encoder layer into the projector without training breaks the answer (0.62 to 0.49 overall): the projector was trained on the final layer's output and does not accept another layer's without retraining. A checkpoint built with the VoxParadox recipe \cite{voxparadox}, prompt-conditioned layer mixing with preference tuning, achieves 0.37 on the conflict pairs against 0.35 for its base model.

\subsection{Supervised Fine-tuning}  
%Before training anything we tried seven interventions that leave the weights frozen: twelve prompt wordings, few-shot audio examples, activation steering, subtracting the transcript-only logits, learned layer mixing, feeding the probe's best layer into the untrained projector, and the recipe of \cite{voxparadox}. None lifts the conflict answer above 0.56, and each one that moves it does so by losing part of the agreement arm; full numbers are on the project page\footnote{\url{https://github.com/parw8649/audio-llm-utility-gap}}. 
Our probing results demonstrated the representation utility gap, \textit{i.e.}, LLM does not use the information available in audio features. Hence, we retrain the projector between the audio encoder and the language model to make paralinguistic information more salient to the LLM. After fine-tuning the projector, AUC improves on every dataset except MDVR-KCL, which has 37 speakers, and E-DAIC, where it drops (Table~\ref{tab:main}). On the Pitt conflict segments the fine-tuned answer stays at 0.52, while the agreement segments rise to 0.86. On E-DAIC the retrained projector lands near the probe range, consistent with a task the words carry and the voice does not. A projector trained on PC-GITA lifts the NeuroVoz answer from 0.70 to 0.74. One trained on NeuroVoz lifts PC-GITA from 0.56 to 0.68. On Parkinson's and Alzheimer's the deficit sits in the readout rather than the representation, since the probe reads the signal at every depth and retraining the projector recovers it. The direction test and the fine tuning are on Qwen2.5-Omni, and we release the segment lists, scripts and every clip's score\footnote{\url{https://github.com/parw8649/audio-llm-utility-gap}}.


%\subsection{Discussion}\label{sec:discussion}

%Read together, the results have one shape. On Parkinson's and Alzheimer's the clinical signal is present at every depth of the model, the answer does not use it, and on all three conditions what the answer uses instead is the words. Every training-free fix that moves the conflict answer pays for it elsewhere, and the best prompt wording moves Qwen2-Audio only from 0.35 to 0.48; the two fixes that train one small component recover most of the Alzheimer's gap (a readout on the frozen states 0.78, the retrained projector's own answer 0.80), and the projector fix moves the MDVR-KCL answers the same way. The deficit therefore lives in the readout, not in the representation: these models are better clinical encoders than clinical reporters, and evaluations that score their answers only where words and diagnosis agree cannot tell a model that listens from one that reads.
\begin{comment}
\noindent\textbf{Limitations} The layer-wise probing and proposed solutions are  are on one model; the nine-model
evidence is the answer-side conflict test. MDVR-KCL has 37 speakers and its projector gain is
not significant on its own; its before number uses the voice wording, the retrained projector the
recording wording. ADReSSo transcripts are ASR output (whisper-large-v3) and its conflict
group has 19 segments; the ADReSS-2020 test labels were recovered by matching transcripts to Pitt
(107 of 108 training labels matched, none wrong), pending the official labels. The E-DAIC 30\,s clips were not filtered for
interviewer speech and will be re-cut. On E-DAIC, scores track PTSD severity
as closely as depression severity (0.60 against 0.59). We release segment identifiers, scripts and every clip's score.
\end{comment}

\section{Conclusions}  
In this paper, we showed how Audio LLM models could under-perform in diagnosing three different clinical conditions. We then inspected the audio representations through probing to learn whether the information necessary for correct diagnosis is present. The probing results demonstrated a wide utility gap, \textit{i.e.}, the LLM does not effectively use the information available from the audio encoder, hence, resulting in lower performance. Similar to previous work \cite{voxparadox}, we found a strong language bias in Audio LLMs as they are more likely to follow the spoken words rather than paralinguistic speech. To alleviate this problem and narrow the utility gap, we performed a supervised fine-tuning limited to the projector layer keeping the rest of the model's weights frozen {\color{red} resulting in a significant performance gain in AD and MDD detection.} This work demonstrated that the hidden state representations carry necessary paralinguistic information in their representation and future work should focus on further narrowing the utility gap.

\section*{Acknowledgments}
Research was sponsored by the ARO and was accomplished under Cooperative Agreement Number W911NF-25-2-0040. Work was also supported by the NIMH under Award Number R61MH135407. The views and conclusions contained in this document are those of the authors and not of the  U.S. Government. The Pitt corpus was supported by NIA grants AG03705 and AG05133. The authors declare no conflict of interest.
\section*{Compliance with Ethical Standards}
This study was a secondary analysis of existing recordings only acquired under each distributor's data agreement. No new recordings were made, no audio is redistributed, and no ethical approval was required for this secondary analysis.
% PI to confirm before submission: whether USC IRB review was required for this secondary analysis,
% and if so the protocol number; ICASSP asks for the statement either way.


\let\oldbibliography\thebibliography \renewcommand{\thebibliography}[1]{\oldbibliography{#1}\setlength{\itemsep}{0pt}\setlength{\parsep}{0pt}\setlength{\parskip}{0pt}\raggedright} \sloppy \setlength{\emergencystretch}{1em} \bibliographystyle{IEEEbib}
\bibliography{refs}
\begin{comment}
    
\clearpage
\begin{table*}[!t]
\centering\small
\caption{Every measurement in one place. Internal review copy, cut before submission.}
\begin{tabular}{@{}p{0.66\textwidth}p{0.30\textwidth}@{}}
\toprule
measurement & AUC \\
\midrule
encoder probe, nested: PC-GITA / Neurovoz / KCL read / Pitt & 0.92 / 0.96 / 0.89 / 0.78 \\
encoder probe, curve peak, Pitt / language-model stages, Pitt & 0.82 / 0.86 \\
model's own answer: PC-GITA / Neurovoz / KCL / E-DAIC / Pitt & 0.54 / 0.55 / 0.59 / 0.65 / 0.62 \\
answer-position probe, Pitt: final layer / best layer / conflict & 0.80 / 0.83 / 0.69 \\
transcript only, Pitt (clean text) / ADReSSo (whisper-large-v3) & 0.68 / 0.68 \\
E-DAIC transcript vs own audio answer: $r$ / identical Yes-No & 0.83 / 96\% \\
words only, Pitt: Qwen2-Audio / omni / diva / mimo & 0.72 / 0.72 / 0.71 / 0.57 \\
words check, Pitt: within agreement / on conflict & 0.95 / $\le$0.58 \\
words check, ADReSSo: within agreement / on conflict (n=19) & 0.99 / 0.39 to 0.61 \\
DAIC conflict, nine models: 0.36 or lower / below 0.5 & 6 of 9 / 7 of 9 \\
Pitt conflict, qwen2.5-omni (No on every clip): conflict / agreement & 0.44 / 0.77 \\
Pitt conflict, Qwen2-Audio: conflict / agreement & 0.41 / 0.70 \\
fixes on conflict: prompts (omni, Qwen2-Audio) / subtraction / steering / few-shot / plug & 0.34, 0.48 / 0.50 / 0.54 / 0.46 / 0.55 \\
text subtraction at $\alpha = 2.0$: conflict / agreement & 0.67 / 0.18 \\
VoxParadox-recipe checkpoint on conflict: checkpoint / base & 0.37 / 0.35 \\
layer mixer vs one nested layer / prompt-conditioned mixer gain & 0.909 vs 0.916 / +0.0004 \\
trained readout, Pitt: overall / agreement / conflict & 0.78 / 0.86 / 0.61 \\
projector retrained, Pitt: overall / agreement / conflict & 0.80 / 0.89 / 0.57 \\
projector retrained on its default final layer (no plug), Pitt: overall / agreement / conflict & 0.76 / 0.86 / 0.48 \\
layer 24 minus final-layer projector (paired): overall / conflict & +0.035 [$-$0.004, 0.073] / +0.089 [0.004, 0.177] \\
projector retrained, KCL: before / after & 0.65 / 0.75 \\
paired gains, Pitt (95\% intervals): projector / readout & +0.18 [0.11, 0.24] / +0.17 [0.10, 0.23] \\
paired gains on the conflict arm, Pitt: readout / projector & +0.19 [0.04, 0.34] / +0.16 [0.02, 0.30] \\
Pitt steering at the probe's layer 22, $c=1$: conflict / agreement (random direction) & 0.56 / 0.42 (0.58 / 0.54) \\
audio answer minus clean transcript, Pitt (paired) & $-0.06$ [$-0.12$, $+0.01$] \\
ADReSSo: probe nested / probe peak / answer / transcript only & 0.87 / 0.89 / 0.66 / 0.68 \\
ADReSS-2020 official split: probe AUC / accuracy / challenge acoustic, linguistic baselines & 0.80 / 73\% / 62.5\%, 75\% \\
ADReSS-2020 answer: test 48 / train 108 / pooled 156 & 0.78 / 0.62 / 0.68 \\
ADReSS-2020 pooled probe: nested / peak & 0.82 / 0.87 \\
cross dataset, Parkinson's: PC-GITA to KCL / Neurovoz to KCL & 0.86 / 0.74 \\
cross dataset, four further models: granite / voxtral / omni / diva & 0.81 / 0.85 / 0.83 / 0.84 \\
language-model stages on KCL read: min / max & 0.60 / 0.69 \\
classical acoustic floor (eGeMAPS): Pitt / ADReSSo / ADReSS-2020 / KCL & 0.62 / 0.73 / 0.70 / 0.80 \\
recording-only floor: Italian dataset (excluded) / other datasets / Pitt six measurements, duration alone & 0.995 / 0.45 to 0.80 / 0.65, 0.61 \\
Pitt age: age alone AUC / 64 matched pairs, 268 segments: probe / answer / age alone & 0.70 / 0.79 / 0.66 / 0.38 \\
matched 300\,s window, depression: flamingo 3 / omni / kimi / omni at 30\,s & 0.86 / 0.84 / 0.83 / 0.67 \\
\bottomrule
\end{tabular}
\end{table*}
\end{comment}

\end{document}
